\documentclass[10pt,a4paper]{amsart}
\usepackage[margin=19mm]{geometry}
\usepackage{amsmath,amssymb,mathtools}
\usepackage[scr=rsfs,cal=euler]{mathalfa}
\usepackage{xfrac}
\usepackage[unicode,hidelinks,bookmarks=false]{hyperref}
\newcommand{\ie}{i.e.\ }
\newcommand{\cf}{cf.\ }
\newcommand{\resp}{respectively\ }
\newcommand{\Subsection}{\subsection}
\providecommand{\nobreakdash}{\nobreak\hbox{-}}
\setlength{\emergencystretch}{2em}
\multlinegap0pt
\usepackage{bm}
\usepackage{bbm}
\usepackage{dsfont}
\usepackage{xspace}
\usepackage{cancel}
\usepackage{mathtools}
\usepackage{amsthm}
\makeatletter
\@ifundefined{theorem}{\newtheorem{theorem}{Theorem}}{}
\@ifundefined{lemma}{\newtheorem{lemma}[theorem]{Lemma}}{}
\makeatother
\newcommand{\N}{\mathbb{N}}
\newcommand{\floor}[1]{\lfloor#1\rfloor}
\newcommand{\norm}[1]{\left\lVert#1\right\rVert}
\title[More stability and convergence results for higher-order Wien]{More stability and convergence results for higher-order Wiener-Wintner systems}
\author{Jacob Folks}
\date{Source: arXiv:2510.19082v1 (CC BY 4.0)}
\begin{document}
\maketitle

\noindent\emph{Source and licence.} More stability and convergence results for higher-order Wiener-Wintner systems. Version arXiv:2510.19082v1 (CC BY 4.0), distributed under the Creative Commons Attribution 4.0 International licence, as recorded in the arXiv licence metadata for this exact version and shown on the abstract page https://arxiv.org/abs/2510.19082v1. Full rights notice: licenses/arxiv-2510.19082-CC-BY-4.0.txt.\par\smallskip

\noindent\textit{Editorial scope.} Counted unit: the Theorem whose source label is \texttt{T:polynomial return times} in the article (source0.tex lines 972--978, source label \texttt{T:polynomial return times}), delivered together with the article's own complete original proof of it (source0.tex lines 980--1027, one \texttt{proof} environment of 48 lines). No mathematics is written, repaired or re-derived: statement and proof are the article's own text line for line. Required context retained uncounted: L:maxIneq (lines 367--370); T:polynomial WW estimate (lines 916--922). Every definition, display and label of those lines is retained unchanged. Omitted from this package and not redistributed: 1--366, 371--915, 923--971, 979--979, 1028--1603. Nothing inside a retained excerpt is omitted or altered: every retained body is delivered exactly as the article's lines are written, so no omission receipt of any kind accompanies this package. Citations: the retained excerpts carry no citation command at all, so no bibliography entry of the article is transcribed and no reference is redistributed. Reference adaptation: none was needed and none was made; every reference inside the delivered unit resolves inside it, so no unresolved reference or citation is produced. Printed slips kept literal: no printed word, symbol, hypothesis or number of the counted statement or of its proof was altered. Layout and class: the article's own class is \texttt{amsart} with packages including biblatex, geometry, graphicx, babel, amsthm, bbm, mathtools, hyperref, ifthen, tikz, tikz-cd, blindtext, scrextend, amssymb; the wrapper is the lane's pinned \texttt{amsart} editorial wrapper at 10pt on A4, which restates the theorem-like environments the retained text needs and adds the article's own macro definitions that the retained text needs (\texttt{\textbackslash N}, \texttt{\textbackslash floor}, \texttt{\textbackslash norm}), with extra wrapper packages (bm, bbm, dsfont, xspace, cancel, mathtools). The delivered numbering is the unit's own. Assets: no retained span contains a figure environment, a graphics-inclusion command, a PDF-page inclusion, a table or a TikZ picture; no image file is copied into this package, the package declares no asset dependency and no image file is redistributed with it. Licence: version arXiv:2510.19082v1 of this article is distributed under the Creative Commons Attribution 4.0 International licence, as recorded in the arXiv licence metadata for this exact version and shown on the abstract page https://arxiv.org/abs/2510.19082v1. A full rights notice is delivered at licenses/arxiv-2510.19082-CC-BY-4.0.txt. Characters: every retained excerpt is pure ASCII, so the delivered engine, XeLaTeX with its default Latin Modern text font, reports no missing character.
\begin{lemma}[Maximal inequality]\label{L:maxIneq}
Let $(X, \mathcal{F}, \mu, T)$ be a measure-preserving system, and $p \in (1, \infty)$. For every real-valued function $f \in L^p(\mu)$, we have
\[ \norm{  \sup_N \frac{1}{N} \sum_{n=1}^N f \circ T^n}_p \leq \frac{p}{p-1} \norm{f}_p \, . \]
\end{lemma}

\begin{theorem} \label{T:polynomial WW estimate}
Let $J \in \N$ and $k\in \N \cup \{0\}$, and $a_1, \dots, a_J\in \mathbb{Z}$ be distinct and nonzero. Then there exists a constants $C_{J, k, a}$ and $N_{J, a}$ such that for all invertible dynamical system $(X, \mathcal{F}, \mu, T)$ and $f_1, \dots, f_J\in L^\infty (\mu)$, all bounded by 1, we have
\begin{align}\label{E:poly bound}
\left\Vert \sup_{t_1, \dots, t_k} \left|\frac{1}{N}\sum_{n=1}^N e^{2\pi i p_{k, t}(n)} \prod_{j=1}^J f_j \circ T^{a_j n} \right| \right\Vert_2 &\leq C_{J, k, a}\left(\frac{1}{\floor{\sqrt{N}}^{1/2^{k+J-2}}} + \left[W_N^{k+J-1}(f_1)\right]^{1/2^{k+J-2}}\right)
\end{align}
for all $N> N_{J, a}$, where $p_{k, t}(n) = t_1 n + t_2n^2 + \dots t_kn^k$ is the polynomial with coefficients $t_i$.
\end{theorem}

\begin{theorem}\label{T:polynomial return times}
Let $(X, \mathcal{F}, \mu, T)$ be an invertible dynamical system, and let $f_1 \in L^\infty(\mu)$ be an $L^2$ limit of $k+J-1$-th order WW functions. Then for any functions $f_2, \dots, f_J\in L^\infty(\mu)$ all bounded by 1, there exists a set $X_{f_1, \dots, f_J, k}$ of full measure such that for any $x\in X_{f_1, \dots, f_J, k}$ and $a_1, \dots, a_J$ all distinct and nonzero and any dynamical system $(Y, \mathcal{G},\nu, S)$ and $g\in L^\infty(\nu)$ and any integer valued polynomial $P$ of degree less than or equal to $k$, the averages
\begin{align*}
\frac{1}{N}\sum_{n=1}^N g(S^{P(n)} y) \prod_{j=1}^J f_{j}(T^{a_j n} x)
\end{align*}
converge to zero $\nu$-a.e.
\end{theorem}

\begin{proof}
Since there are countably many possibilities for $a_1, \dots, a_J$, it follows that we may fix them and intersect over a countable collection of full measure sets at the end.

Consider first the case in which $f_1$ is itself a $k+J-1$-th order WW function of power type $\alpha$. Using the previous Theorem \ref{T:polynomial WW estimate}, it follows for sufficiently large $N$ that we have
\begin{align*}
\left\Vert \sup_{t_1, \dots, t_k} \left|\frac{1}{N}\sum_{n=1}^N e^{2\pi i p_{k, t}(n)} \prod_{j=1}^J f_j \circ T^{a_j n} \right| \right\Vert_2 &\leq C_k\left(\frac{1}{\floor{\sqrt{N}}^{1/2^{k+J-2}}} + \left[\frac{C_{f_1}}{N^\alpha}\right]^{1/2^{k+J-2}}\right) \leq \frac{C'}{N^{\alpha / 2^{k + J - 2}}}
\end{align*}
after consolidating constants. Picking $\gamma \in \N$ to satisfy $\alpha \gamma > 2^{k + J}$, it follows that 
\begin{align*}
\int \sup_{t_1, \dots, t_k} \left|\frac{1}{\floor{N^\gamma}}\sum_{n=1}^{\floor{N^\gamma}} e^{2\pi i p_{k, t}(n)} \prod_{j=1}^J f_j (T^{a_j n} x) \right|^2 d\mu(x) &\leq \frac{C'}{\floor{N^\gamma}^{\alpha/2^{k+J-1}}} \leq \frac{C'}{N^2}
\end{align*}
for sufficiently large $N$, where the constant is allowed to change, but picks up no $N$ dependence. These terms are summable in $N$, and by the monotone convergence theorem, we have
$$\sum_{N=1}^\infty \sup_{t_1, \dots, t_k} \left|\frac{1}{\floor{N^\gamma}}\sum_{n=1}^{\floor{N^\gamma}} e^{2\pi i p_{k, t}(n)} \prod_{j=1}^J f_j (T^{a_j n} x) \right|^2 < \infty $$
for all $x$ in a set of full measure $X_{f_1, \dots, f_{J}, k}$.

For any $x\in X_{f_1, \dots, f_J, k}$, applying the spectral theorem for any $g\in L^\infty(\nu)$, we have
\begin{align*}
\left\Vert \frac{1}{\floor{N^\gamma}} \sum_{n=1}^{\floor{N^\gamma}} g \circ S^{P(n)} \cdot \prod_{j=1}^J f_j (T^{a_j n} x)\right\Vert_{L^2(\nu)}^2 &= \int_0^1 \left| \frac{1}{\floor{N^\gamma}} \sum_{n=1}^{\floor{N^\gamma}} e^{2\pi i P(n) t} \prod_{j=1}^J f_j (T^{a_j n} x) \right|^2 \, d\sigma_g(t)  \\
& \leq \Vert g \Vert_2^2 \cdot \sup_t \left| \frac{1}{\floor{N^\gamma}} \sum_{n=1}^{\floor{N^\gamma}} e^{2\pi i P(n) t} \prod_{j=1}^J f_j (T^{a_j n} x) \right|^2 \\
& \leq \Vert g \Vert_2^2 \cdot \sup_{t_1, \dots, t_{k}} \left| \frac{1}{\floor{N^\gamma}} \sum_{n=1}^{\floor{N^\gamma}} e^{2\pi i p_{k, t}(n)} \prod_{j=1}^J f_j (T^{a_j n} x) \right|^2
\end{align*}
and again these terms are summable. By the monotone convergence theorem in $\nu$, it follows that the integrand
\begin{align}\label{F:average along}
\left|\frac{1}{\floor{N^\gamma}} \sum_{n=1}^{\floor{N^\gamma}} g( S^{P(n)} y) \prod_{j=1}^J f_j (T^{a_j n} x) \right|^2
\end{align}
convergence to zero $\nu$-a.e. Since the summands are bounded, we may extend the convergence from the subsequence $\floor{N^\gamma}$ to $N$ by the standard argument. We outline this argument here, for later comparison: Denote the summands in $n$ of the average (\ref{F:average along}) as $a_n$, and observe that $|a_n| \leq 1$. Moreover, let us ignore the square root. Let $N\in \mathbb{N}$ be arbitrary, and pick $M\in \mathbb{N}$ to satisfy
$$M^\gamma \leq N \leq (M+1)^\gamma \, .$$
We observe that
\begin{align*}
\left| \frac{1}{N}\sum_{n=1}^N a_n\right| &\leq \left|\frac{1}{\floor{M^\gamma}}\sum_{n=1}^{\floor{M^\gamma}} a_n\right| + \frac{1}{M^\gamma} \sum_{n = \floor{M^\gamma}+1}^{\floor{(M + 1)^\gamma}} |a_n| \leq \left|\frac{1}{\floor{M^\gamma}}\sum_{n=1}^{\floor{M^\gamma}} a_n\right| + \frac{\floor{(M + 1)^\gamma}-\floor{M^\gamma}}{M^\gamma} 
\end{align*}
As $N\to \infty$, it follows that $M \to \infty$ also: the averages along $\floor{M^\gamma}$ goes to zero by assumption, and the remainder term is $O(M^{-1})$, and also goes to zero. Hence, we obtain convergence along the entire subsequence.

Now, suppose that $F_m$ are $j+K-1$-th order WW functions of power type converging to $f_1$. Without loss of generality, take $\Vert F_m - f_1 \Vert_2< \frac{1}{m}$ for all $m\in \N$. Using the maximal ergodic theorem (here referenced as lemma \ref{L:maxIneq}), we observe for each $m\in \N$ and any $g \in L^\infty(\mu)$ that
\begin{align*}
&\left\Vert \frac{1}{N}\sum_{n=1}^N g\circ S^{P(n)} \cdot \prod_{j=1}^J f_j\circ T^{a_j n} - \frac{1}{N}\sum_{n=1}^N g\circ S^{P(n)} \cdot F_m \circ T^{a_1n} \cdot \prod_{j=2}^J f_j\circ T^{a_j n} \right\Vert_2 \\
&\leq \left\Vert \frac{1}{N}\sum_{n=1}^N \left| f_1 - f_M \right| \circ T^{a_1 n}  \right\Vert_2 \leq 2\Vert f_1 - F_m \Vert_2 \leq \frac{2}{m^2}
\end{align*}
and these terms are summable. By the monotone convergence theorem, it follows that the integrand must be converging to zero almost everywhere as $m\to \infty$.

Hence, for any $x\in \cap_{m=1}^\infty X_{F_m, f_2, \dots, f_J, k}$, and for any particular value of $m$, we see
\begin{align*}
&\limsup_N \left| \frac{1}{N}\sum_{n=1}^N g\circ S^{P(n)} \cdot \prod_{j=1}^J f_j\circ T^{a_j n}\right| \\
&\leq \limsup_N \left| \frac{1}{N}\sum_{n=1}^N g\circ S^{P(n)} \cdot \prod_{j=1}^J f_j\circ T^{a_j n} - \frac{1}{N}\sum_{n=1}^N g\circ S^{P(n)} \cdot F_m \circ T^{a_1n} \cdot \prod_{j=2}^J f_j\circ T^{a_j n} \right| + 0 \\
&\leq \sup_N \left| \frac{1}{N}\sum_{n=1}^N g\circ S^{P(n)} \cdot \prod_{j=1}^J f_j\circ T^{a_j n} - \frac{1}{N}\sum_{n=1}^N g\circ S^{P(n)} \cdot F_m \circ T^{a_1n} \cdot \prod_{j=2}^J f_j\circ T^{a_j n} \right| 
\end{align*}
which converges to zero as we let $m\to \infty$.
\end{proof}

\end{document}
